% Related Work / Literature Review
\subsection{Overview}
This section reviews prior work in human motion tracking, sensor-based and vision-based techniques, AI model generation frameworks, edge AI deployment, and preprocessing strategies, situating the proposed framework within existing research.

\subsection{Application Domains of Motion Tracking}
Human motion tracking is a foundational technology across healthcare, sports and fitness, virtual/augmented reality (VR/AR), robotics, and emerging interactive systems. In healthcare, wearable and ambient sensors enable gait analysis, fall detection, and rehabilitation monitoring, supporting personalized interventions \cite{article:human_motion_tracking_survey}. Sports applications leverage inertial and biomechanical data for performance optimization and injury prevention, while VR/AR systems depend on accurate skeletal and positional tracking for immersive user experiences. Robotics integrates motion tracking for navigation, manipulation, and human-robot interaction. These domains collectively motivate the need for accessible, adaptable, and efficient AI-driven motion tracking frameworks.

\subsection{Sensor-Based Motion Tracking}
Inertial Measurement Units (IMUs)—combining accelerometers and gyroscopes—are widely adopted for motion tracking due to their low power consumption, cost efficiency, and suitability for edge deployment. Survey work focusing on upper-limb tracking \cite{article:motion_tracking_IMU_upper_limb} highlights challenges including sensor drift, cumulative error, calibration, and biomechanical constraints. Multi-sensor fusion approaches (e.g., IMU + magnetometer + pressure sensors) can improve robustness but increase complexity. The proposed framework emphasizes reliable single-IMU pipelines with extensibility for multi-sensor integration.

\subsection{Vision-Based Motion Tracking}
Computer vision approaches using stereoscopic cameras, markerless tracking, and pose estimation provide rich spatial representations of human motion \cite{article:motion_tracking_using_vision}. While vision systems enable full-body reconstruction and are used in animation and clinical analysis, they face challenges in occlusion, environmental variability, privacy, and power consumption. Recent advances (e.g., MediaPipe pose and hand tracking \cite{web:Google_Mediapipe}) have improved accessibility but remain less feasible for ultra-low-power wearable devices compared to IMU-based solutions.

\subsection{AI in Motion Tracking}
AI models transform raw motion data into semantic interpretations—classifying activities, identifying postures, and forecasting trajectories. Traditional machine learning approaches (e.g., Support Vector Machines (SVM) \cite{cortes1995svm}, Random Forests \cite{breiman2001random}) have been widely used due to interpretability and compact runtime footprints suitable for edge devices. Deep learning, including CNNs, RNNs/LSTMs, and emerging transformer architectures, automates hierarchical feature extraction and improves recognition performance, especially on complex multi-modal datasets. However, deep learning models frequently demand larger annotated datasets, higher memory usage, and non-trivial quantization strategies during edge deployment. The proposed framework strategically focuses on classical models and lightweight neural architectures to balance performance and deployability.

\subsection{Model Generation Frameworks}
Several platforms have emerged to simplify AI model creation:
\begin{itemize}
    \item \textbf{Edge Impulse} \cite{web:Edge_Impulse}: Provides end-to-end pipelines for data ingestion, preprocessing, feature generation, model training, and deployment for microcontrollers. It excels in automation but restricts transparency and deep customization in preprocessing logic.
    \item \textbf{MediaPipe} \cite{web:Google_Mediapipe}: Offers modular perception pipelines (pose, hand, object tracking) with optimized graph execution. Primarily vision-oriented, less suited for bespoke sensor-driven workflows requiring custom segmentation strategies.
    \item \textbf{Teachable Machine} \cite{web:Google_TeachableMachine}: Enables rapid prototyping of classifier models through a no-code interface, but lacks instrumentation for advanced feature engineering and hardware-aware deployment.
    \item \textbf{TensorFlow Lite / TinyML toolchains} \cite{web:TensorFlowLite, warden2019tinyml}: Provide model conversion, quantization, and runtime inference optimizations for edge devices. These toolchains assume a separate model development environment and offer limited UI support for high-fidelity preprocessing.
\end{itemize}
Compared to these platforms, the proposed framework contributes a novel interactive preprocessing workflow with draggable time window segmentation and direct per-activity window refinement—features absent from current commercial ecosystems.

\subsection{Data Preprocessing and Windowing}
Accurate activity recognition depends heavily on the quality of segmentation and windowing strategies. Studies such as \cite{banos2014wisdm} examine window size impacts on activity recognition performance; shorter windows increase temporal resolution but may reduce feature stability, while longer windows capture more context at the cost of responsiveness. Conventional preprocessing pipelines require manual code modification to adjust filtering, smoothing, or artifact removal. The proposed framework advances usability by enabling direct visual alignment of raw and preprocessed signals, interactive window adjustments, and on-the-fly re-labeling—lowering cognitive and technical barriers for iteration.

\subsection{Feature Extraction and Representation}
Handcrafted time-domain (mean, variance, zero-crossing rate) and frequency-domain features (spectral energy, dominant frequency) remain effective for classical ML models when extracted systematically \cite{oppenheim1999discrete}. Feature selection can improve memory footprint and inference speed on microcontrollers. While end-to-end deep networks can learn hierarchical features implicitly, explicit feature extraction improves interpretability and deployment transparency in regulated or safety-critical domains. The framework incorporates pluggable feature computation modules and supports per-window diagnostics to guide refinement.

\subsection{Edge Deployment and Optimization}
Resource constraints define edge deployment challenges—limited RAM/flash, absence of floating-point units, and power budgets. Optimization strategies include model pruning, quantization, fixed-point arithmetic, and algorithmic simplification \cite{banbury2021mlperf}. Classical models like Random Forests (with constrained depth and tree count) and SVMs (with reduced support vector sets) remain strong candidates for IMU-based HAR tasks. The framework’s code generation component performs platform-tailored export, emitting structured C/C++ code with constraint-aware parameter transpilation. Multi-objective optimization modes (accuracy, speed, power, balanced) enable decision-making aligned with deployment goals.

\subsection{Gaps in Existing Work}
Despite mature tooling for edge AI development, key gaps persist:
\begin{enumerate}
    \item \textbf{Limited Interactive Preprocessing}: Few systems allow users to visually inspect raw vs. filtered sensor signals and manually refine segmentation inside the same environment.
    \item \textbf{Opaque Feature Pipelines}: Automated feature generators hide derivation steps and hinder educational uses or research replication.
    \item \textbf{Non-Extensible Deployment Export}: Code generation output is often monolithic, complicating custom optimization passes or hardware adaptation.
    \item \textbf{Accessibility vs. Control Trade-off}: Platforms favor either simplicity (Teachable Machine) or configurability (TensorFlow), rarely combining both for motion sensor workflows.
\end{enumerate}
The proposed framework addresses these gaps by unifying interactive preprocessing, transparent feature extraction, modular training, and optimization-aware code generation under an extensible architecture.

\subsection{Summary}
Prior research and industry platforms provide strong foundations in motion tracking, yet none deliver an integrated, transparent, and user-friendly workflow optimized for IMU-driven edge deployment with interactive segmentation. This work builds upon established algorithms and edge ML practices while introducing novel contributions that democratize development, accelerate iteration, and improve deployment readiness for real-world motion tracking applications.
